% One-page presenter's card for the 16 Oct meeting. Compile: tectonic presenter_card.tex
% Numbers: docs/NEXT_STEPS.md section 1a and DECISIONS_LOG.md #95-107.
% AI assistance: Claude drafted this card from the project's results.
\documentclass[10pt]{article}
\usepackage[a4paper,margin=1.3cm]{geometry}
\usepackage{multicol,enumitem,booktabs}
\setlist{nosep,leftmargin=1.2em}
\setlength{\parindent}{0pt}
\setlength{\columnsep}{1.2em}
\pagestyle{empty}
\newcommand{\h}[1]{\vspace{4pt}{\large\bfseries #1}\par\vspace{2pt}}

\begin{document}
{\LARGE\bfseries Friday 16 Oct: presenter's card}\hfill{\small slides: \texttt{meeting\_plots/oct16/slides.pdf}}

\vspace{4pt}
\textbf{Open with:} ``T1 is complete in 17 regions. I built the T3 test
harness and finished T2 tier 1. The headline: terrain from elevation data
predicts hiking speed in regions the model has never seen; OSM trail tags
barely do. It survived every check I could think of.''

\begin{multicols}{2}
\h{Numbers to know}
\begin{itemize}
\item \textbf{Data:} 17 regions, 2,876 hiking traces; walking set 1,992
  outings, 290,640 segments of 50 m, 4,702 h.
\item \textbf{Tobler:} 13\% worse than our curve. Real hikers peak near
  \textbf{4.2 km/h}, not 6.
\item \textbf{Unseen region} (skill vs slope curve): terrain
  \textbf{+2.1\%} (trees +2.8\%), tags +0.9\% (+0.4\%). Terrain beats
  tags in \textbf{16 of 17} regions.
\item \textbf{Harder tests:} a whole zone left out +1.9\%; West
  $\leftrightarrow$ East +2.6\%; newer hikes +2.2\%; same at every GPS
  rate.
\item \textbf{Whole hike:} time error 13.6\% $\to$ 9.6\% with the hiker's
  own pace (first km) + terrain.
\item \textbf{Same hiker, same trail:} down is 8\% faster than up on gentle
  grades, 22\% on steep ones. Tobler says 42\%.
\item \textbf{Modes:} runners = walkers $\times$2.5; bikes need their own
  curve.
\end{itemize}

\h{New since last week}
\begin{itemize}
\item Why tags fail: not different meanings (within a hike, harder
  sac\_scale is slower in every region). They're missing on much of the
  trail (1--96\% by region).
\item Sparse GPS traces read up to 11\% slow: how we measure, not slower
  hikers.
\item A 2026 paper (Huang et al., Hong Kong) uses boosted trees with
  weather, one city, random split: no transfer test.
\item The 5-number formula travels less well than the full model (+0.6\%
  for an unseen zone). Say so.
\end{itemize}

\columnbreak
\h{Ask Danni}
\begin{enumerate}
\item Framing: ``terrain transfers, tags don't''. Agree?
\item Embeddings: keep as ``helps locally, transfers less''?
\item Lidar understory: $\sim$12 GB for the White Mountains. Do it?
\item Data release: license (ODbL) and lab policy.
\item Cluster account: \texttt{engr-class-any} or the lab's?
\item Fix sparse traces in a T1 v2, or just report it?
\item Which 3--4 figures to keep (the draft is at 5 pages).
\end{enumerate}
\textit{Same order as slide 14.}

\h{Hard questions, short answers}
\begin{description}[leftmargin=0pt,style=nextline]
\item[``Isn't +2\% small?''] Per 50 m, yes: GPS noise and the hiker
  dominate. But it's consistent (17 of 17 regions in direction), its
  interval excludes zero, and it's more than twice what tags give.
\item[``Is terrain just a stand-in for the region?''] No: it adds
  $\sim$1.9 points over each region's own average, and holds with a whole
  zone left out.
\item[``Why no deep learning or images?''] About 2,000 outings is small.
  We tested frozen satellite embeddings: +1.1\% in new regions, less
  reliable than terrain.
\item[``What about weather?''] Date and time of day add nothing in new
  regions (seasons differ by region). Weather itself isn't tested yet.
\item[``What did the AI do?''] Claude helped write and run the code and
  draft the documents; I review the work and I'm responsible for it.
\end{description}
\end{multicols}
\end{document}
